% Source: source/普通高考/2014/2014福建理.pdf, page 1, question 5.
% PDF embedded-bitmap attempt: pdfimages -list found no image objects; source chart is vector line art.
% Type: program flowchart (schematic topology and labels are the mathematical content).
% Node -> directed-edge list: 开始 -> S=0,n=1 -> S=S+2^n+n -> n=n+1 -> S>=15;
% 是 -> 输出 S -> 结束; 否 -> right loop -> S=S+2^n+n.
% solid segments: all node borders and directed connectors; dashed segments: none.
% The return arrow points left at the connector above the update box, followed by one
% downward arrow into that box; all node text is locally zero-indented and centered.
\begin{tikzpicture}[
    x=0.90cm,
    y=1.12cm,
    >=Stealth,
    line width=0.45pt,
    line cap=butt,
    line join=miter,
    flowtext/.style={anchor=center,align=center,
        execute at begin node={\setlength{\parindent}{0pt}}},
    every node/.style={flowtext,font=\normalsize},
    flowbox/.style={draw=black,flowtext,
        inner xsep=4pt,inner ysep=2.5pt},
    terminal/.style={flowbox,rounded corners=0.18cm,minimum width=1.45cm,minimum height=0.68cm},
    process/.style={flowbox,minimum height=0.68cm},
    decision/.style={flowbox,diamond,aspect=2.8,minimum height=0.96cm},
    io/.style={flowbox,trapezium,trapezium left angle=72,trapezium right angle=108,
        minimum height=0.74cm},
]
    % Tight bounding box avoids an A4-like page around the standalone chart.
    % Keep a small exterior margin around the widest update box so its border
    % is not clipped by the standalone bounding box.
    \path[use as bounding box] (-2.45,-0.35) rectangle (4.05,9.65);

    \node[terminal] (start) at (0,8.95) {开始};
    \node[process,minimum width=3.35cm] (init) at (0,7.55) {$S=0,n=1$};
    \node[process,minimum width=4.10cm] (update) at (0,6.10) {$S=S+2^n+n$};
    \node[process,minimum width=2.35cm] (increment) at (0,4.62) {$n=n+1$};
    \node[decision,minimum width=4.05cm] (test) at (0,3.12) {$S\geq 15?$};
    \node[io,minimum width=2.45cm] (output) at (0,1.50) {输出 $S$};
    \node[terminal] (finish) at (0,0.05) {结束};

    \coordinate (loop-right) at (3.25,3.12);
    \coordinate (loop-right-top) at (3.25,6.72);
    % Keep the left-pointing return arrowhead clear of the degree-3 merge;
    % the short undirected link then joins the single downward arrow.
    \coordinate (loop-tip) at (-0.35,6.72);
    \coordinate (loop-top) at (0,6.72);

    \draw[->] (start.south) -- (init.north);
    \draw (init.south) -- (loop-top);
    \draw[->] (loop-top) -- (update.north);
    \draw[->] (update.south) -- (increment.north);
    \draw[->] (increment.south) -- (test.north);
    \draw[->] (test.south) -- node[right] {是} (output.north);
    \draw[->] (output.south) -- (finish.north);

    % 否 branch returns around the right side and re-enters the update connector.
    \draw (test.east) -- node[above] {否} (loop-right) -- (loop-right-top);
    \draw[->] (loop-right-top) -- (loop-tip);
    \draw (loop-tip) -- (loop-top);
\end{tikzpicture}
